\section{Related work}
\label{sec:related}

\Cref{tab:related} positions this study against the closest prior work along
three axes: the base RL family, whether an analytic (through-simulator)
gradient is used, and whether the chaotic/turbulent regime is addressed.

\begin{table*}[tbp]
\centering
\caption{Positioning against the closest prior work. AG $=$ analytic
(through-simulator) policy gradient; ME $=$ maximum-entropy (SAC-style)
objective; \checkmark\ = present, dash = absent.}
\label{tab:related}
\footnotesize
\setlength{\tabcolsep}{4pt}
\resizebox{\textwidth}{!}{%
\begin{tabular}{|l|l|l|c|c|l|}
\hline
Work & Base & AG & ME & Chaotic & Domain \\
\hline
HydroGym / GPPO~\cite{hydrogym2026} & PPO & \checkmark & -- & one $\Re$ & Kolmogorov flow \\
\hline
GI-PPO~\cite{gippo2023}             & PPO & \checkmark (adapt.) & -- & -- & rigid body \\
\hline
RPO~\cite{rpo2025}                  & PPO & \checkmark & -- & -- & control suites \\
\hline
SHAC~\cite{shac2022}                & AC  & \checkmark (short-$H$) & -- & -- & locomotion \\
\hline
AHAC~\cite{ahac2024}                & AC  & \checkmark (adapt.-$H$) & -- & -- & contact-rich \\
\hline
SAPO~\cite{sapo2025}                & SAC & \checkmark & \checkmark & -- & rigid/deformable \\
\hline
Bucci et al.~\cite{bucci2019}       & DDPG & -- & -- & \checkmark & Kuramoto--Sivashinsky \\
\hline
Font et al.~\cite{font2025}         & PPO & -- & -- & \checkmark & separation bubble \\
\hline
This work                  & PPO, SAC & \checkmark & \checkmark & 45 / 300 & Kolmogorov flow \\
\hline
\end{tabular}%
}

\end{table*}

\subsection{RL for flow control and HydroGym}
Since Rabault et al.~\cite{rabault2019}, model-free RL has repeatedly
succeeded on turbulent flow control: separation-bubble
control~\cite{font2025} reports $\sim$\SI{9}{\percent} vs $\sim$\SI{6.8}{\percent}
reduction over classical periodic forcing; an early study~\cite{turbdrl2020}
establishes DRL's applicability to turbulent conditions; real-time
state-dependent bluff-body wake control~\cite{react2026} achieves 2--4$\times$
the performance of model-based baselines in wind-tunnel experiments; deep RL
has also learned energy-efficient collective swimming gaits by exploiting
shed vortices~\cite{verma2018vortices}. A broad
robotics survey~\cite{drlrobotics2025} documents the wider maturity of RL
for real-world control. HydroGym~\cite{hydrogym2026} unifies these under one
solver-agnostic interface, shows zero-shot policy transfer across
environments, and reports the $\ge$\SI{65}{\percent} sample-efficiency gain
for GPPO over PPO on Kolmogorov flow that this paper examines.
FluidGym~\cite{fluidgym2026} is a related PyTorch-based differentiable
benchmark with PPO/SAC baselines. None of these works systematically
compares analytic-gradient injection across a regular and a chaotic regime.

\subsection{Analytic policy gradients through differentiable simulation}
The single closest prior method to GPPO is GI-PPO~\cite{gippo2023}:
it introduces an ``$\alpha$-policy'' that adaptively weights the analytic
gradient against the PPO gradient, down-weighting the analytic term when it
is detected to be high-variance or biased. Our fixed 50/50 mix is the
non-adaptive special case ($\alpha$ held constant). RPO~\cite{rpo2025}
analyzes the on-policy analytic-gradient setting directly, showing that
reparameterization policy gradients with sample reuse optimize a PPO-style
surrogate via backpropagation-through-time (BPTT) and targeting exactly the
training instability this paper measures at $\Re=300$.
The broader family originated in robotics/control. SHAC~\cite{shac2022}
decouples the analytic gradient from the full episode horizon using
short-horizon rollouts and a smooth critic; our 6-step horizon is a direct
adaptation. AHAC~\cite{ahac2024} replaces the fixed horizon with
one adaptively truncated to avoid stiff dynamics. SAPO~\cite{sapo2025}
carries the idea into a maximum-entropy, off-policy framework, testing
whether entropy regularization smooths the analytic-gradient landscape; our
GSAC is the first application of the SAPO idea to fluid dynamics.

\subsection{Analytic-gradient instability in chaotic systems}
The theoretical basis of our hypothesis is Metz et al.~\cite{metz2021}: in
chaotic dynamical systems, backpropagated gradients explode through a
butterfly-effect mechanism rooted in the system Jacobian spectrum, with
clipping, truncated BPTT, or a switch to zeroth-order methods proposed as
remedies. Parmas et al.~\cite{parmas2018pipps} show the same failure
directly in model-based policy search: gradients through long chains of
nonlinear computation become large and effectively random in direction, and
reparameterization (analytic) gradients suffer from it while
likelihood-ratio (model-free) gradients are comparatively robust --- the
same asymmetry this paper's on-policy result partially echoes
(\cref{sec:results}). The fluid-dynamics counterpart is Boffetta and
Musacchio~\cite{boffetta2017}: the largest Lyapunov exponent (estimated in
practice via the classical time-series method of Wolf et
al.~\cite{wolf1985lyapunov}) of homogeneous-isotropic turbulence grows with
$\Re$ and the predictability horizon shrinks inversely with it. Together
these frame our 500-episode
finding that GPPO and GSAC also eventually reach NaN, only later: instability
is inevitable; what differs is only the delay mechanism. That gradient-free
RL \emph{can} control spatiotemporal chaos is not in question --- Bucci et
al.~\cite{bucci2019} do exactly this on the Kuramoto--Sivashinsky equation
--- the question here is whether \emph{analytic-gradient} injection survives
the chaotic regime.

\subsection{Analytic vs.\ zeroth-order gradients; partial differentiability}
Suh et al.~\cite{suh2022} compare analytic and zeroth-order policy gradients
in systems with stiffness and discontinuities, and via a bias--variance
analysis propose an $\alpha$-graded interpolating estimator, of which our
50/50 mix is roughly the $\alpha=0.5$ special case. A 2026
follow-up~\cite{onoda2026} re-tests the same question with a lighter
hybrid-switching scheme. ABPT~\cite{abpt2025} proposes an amended BPTT scheme
combining 0-step and $N$-step returns for partially differentiable rewards,
a problem class close to HydroGym's mixed control-reward / physics-state
structure.

\subsection{Base algorithms and distributed PDE control}
The gradient-free anchors are PPO~\cite{ppo2017} and Soft
Actor-Critic~\cite{sac2018,sacapps2018}. PPO builds on TRPO~\cite{trpo2015}
and Generalized Advantage Estimation (GAE)~\cite{gae2015}; Soft
Q-learning~\cite{softq2017} is the origin of SAC's maximum-entropy
framework; both algorithms' updates in this study use
Adam~\cite{kingma2015adam}. On the distributed structure of PDE actuation, Peitz et
al.~\cite{peitz2024} exploit translational equivariance to turn a
high-dimensional control problem into a multi-agent problem with identical
uncoupled agents; Kolmogorov flow's low-wavenumber sinusoidal actuators sit
in the same distributed-control setting.
